%%% Constitution for the Open SOurce Club at the Ohio State University
%%% distributed under a creative commons 0 (CC-0)
\documentclass{article}
\title{The Opensource Club Constitution}
\author{Isaac Jones -- Martin Jansche -- Michael Benedict}
\date{20 March, 2000\\ Updated 24 March 2016}
\setcounter{secnumdepth}{0}
\usepackage[normalem]{ulem}
\usepackage{graphicx}
\begin{document}
	\maketitle

	\section{Preamble}

	It is important to recognize the danger that bureaucracy can get in the way of doing cool things.  So long as the requirements of the Ohio State University and the realistic needs of the club are met, efforts should be made to minimize bureaucracy.

	As the club has evolved it has become necessary to elect additional officers in order to aid the further growth of the club. While this is not an attempt to create a bureaucracy, it is necessary in order to delegate duties among members in order to help the club run smoothly. At the time of this writing there are currently three elected positions within the club: President (or Benevolent Dictator), Vice President, and Treasurer. These three positions will continue to exist as long as they are required by the Ohio State University.

	\section{Article I - Name, Purpose and Non-Discrimination Policy}

	\subsection{Section 1 - Name}

	This is the Opensource Club at the Ohio State University.

	\subsection{Section 2 - Purpose}

	Our purpose is to write and advocate free software, and to create a community of excellent programmers. We advocate free software through creation of documentation, providing support, and fighting cluelessness.

	The club should take no ``official'' stand on political issues in debate among the opensource community, should adopt no ``official club license,'' should not create strict guidelines for what kind of software should be developed, should not create strict guidelines for how a software project should be managed (including frequency of updates, release models, code repository management, etc). Doing so would violate the spirit of the Preamble. Individual software projects can be managed however the individuals involved in the project see fit.

	Implicit in this is that the club must allow advocacy of proprietary software written for opensource platforms (for instance) or writing of open source software for a proprietary platform (for instance) if some of its members wish to do so. However, programs written by members of the club, under the flag of the club, and distributed by the club should meet some acceptable definition of open source software.

	\subsection{Section 3 - Non Discrimination Policy}

	In recognition of the importance of welcoming diversity for the sake of creativity, and for the benefit of humanity, this club welcomes all people. This means that our club and it's members will not discriminate against any individual(s) for reasons of age, color, disability, gender identity or expression, national origin, race, religion, sex, sexual orientation, or veteran status.

	\section{Article II - Membership: Qualifications and categories of membership}

	\subsection{Part 1 - Membership Categories and Selection Processes}
	Ohio State University guidelines demand that voting membership be limited to currently enrolled students.  In order to meet this requirement, membership will be divided into voting and non-voting categories.

	\begin{itemize}
		\item Non-voting membership can be obtained simply by attendance of meetings or presence on our mailing list.
		\item Voting membership requires that the individual meet all Ohio State University guidelines for voting membership, and that the individual has attended at least one meeting prior to the vote.
	\end{itemize}
	\subsection{Part 2 - Membership Privileges}

	Accounts on the club computers can be given out to any member who requests an account, provided Ohio State University's rules are followed.  These accounts may be revoked due to inactivity or failure to follow the rules.  Key card access to the club office can be granted on request given the consent of the president, provided that Ohio State University rules are followed.  Access may be revoked due to inactivity or failure to follow the rules.
	When it is necessary to take a vote, voting members are allowed to do so.

	\subsection{Part 3 - Membership Removal}

	Should cause for membership removal be raised, a member can be removed in one of two ways:
	\begin{enumerate}
		\item A majority vote of club officers.
		\item A two thirds majority vote of voting members, providing there is at least a week's notice of the vote.
	\end{enumerate}

	\section{Article III - Officer Positions, Duties, Selection, and Removal}

	\subsection{Part 1 - Officer Positions, Duties, Powers and Limitations.}

	There are three officer positions required by the Ohio State University: The President (or Benevolent Dictator), the Vice President, and the Treasurer.  These three positions will continue to exist as long as they are required by the Ohio State University. Other officer positions may be added and removed as needed, and may be open to election or appointed by the president. The responsibilities of the three aforementioned officers are as follows:

	The President shall have the following responsibilities:

	\begin{itemize}
		\item Managing key card access to the office
		\item Assignment of administrative accounts on the club computers
		\item Securing the meeting location
		\item Announcing meeting location and topic
		\item Representing the club at other functions and in public
		\item Appointment of miscellaneous positions
	\end{itemize}

	The Vice President shall have the following responsibilities:
	\begin{itemize}
		\item Keeping the minutes of the meeting
		\item Carrying out the duties of the President when the President is unable to do so
	\end{itemize}

	The Treasurer shall the following responsibilities:

	\begin{itemize}
		\item Produces and presents a quarterly budget
		\item Purchases pizza and other consumables for meetings
		\item Represents our club at E-Council
		\item Produces and Submits applications for funding
	\end{itemize}

	All of these responsibilities may be delegated to other club members, though they are ultimately the responsibility of the aforementioned officers.  The President will probably be the hardest working member of the club, and therefore needs the power to make unilateral decisions for the sake of saving time.  No officer will necessarily have any control over software projects, since by the definition of opensource software, code forks can happen when they need to, and no one therefore has ultimate power over a software project.

	\subsection{Part 2 - Officer Selection}

	Once a year, elections will take place for the officer positions.  This election must be announced at least one week prior to the vote.  Any voting member who meets The Ohio State University's requirements may run for any of the offices excluding extra positions that are designated as being appointed (see III 1).  The individual who receives the most votes will attain that position for the next academic year, unless removed.  Ties will be broken by a vote of the officers in place before the active election.

	\subsection{Part 3 - Officer Removal}

	Leadership is needed in any club and so officers are required to be
present at eight out of ten official meetings per quarter, while it is also
realized that extenuating circumstances are possible.  In order to eliminate
doubts surrounding an officers excuse for missing a meeting, said officer
must notify at least one of the other officers that they will not be in
attendance at least 24 hours in advance.  If an officer is in violation of
the aforementioned attendance rules, then the club members will take a vote
as to whether or not said officer will be allowed to keep his or her
position. For all other cases of misconduct, the officer will face the same
removal process as any other member of the club.

	\subsection{Part 4 - Replacing Officers}

	If an officer is either removed or resigns during their term, the remaining officers shall replace the missing officer given that all rules set forth by The Ohio State University are followed.

	\section{Article IV - Advisor: Qualification Criteria}

	The advisor should be a person of technical experience. Preferably someone who has been involved in Opensource development. The advisor must grok the goals of the club so that he or she does not get in their way. The advisor exists to provide guidance, mentor-ship, and cluefulness. The advisor's term is one academic year.

	University guidelines demand that the advisor (or failing that, the sub-advisor) ``of student organizations must be full-time members of the University faculty or Administrative \& Professional staff.''

	\section{Article V - Meetings of the Organization}

	It should be recognized that those involved in individual software projects will have meetings separate from the meetings of rest of the club. It is encouraged, however, that all club members be invited to these meetings.

	If none of the officers of the club can be in attendance at a given meeting then said meeting will not take place, as decisions cannot be made without the presiding officers. Members can still hold an independent meeting, but it will not be recognized as an official club meeting.

	\section{Article VI - Method of Amending Constitution: Proposals, notices and voting requirements}

	In the case that someone thinks that they need to alter the constitution or by-laws (if any) for bug fixes, adding features, or compatibility with new hardware, the election must be announced at least a week in advance and the vote must be a 2/3 majority. Voting can be in person or via some electronic medium provided that in either case a reasonable amount of certainty of identity can be secured. Amendments and changes should be taken advisedly and considered for a reasonable amount of time before being implemented.

	It is strictly discouraged that the constitution should be amended frequently. If it is to be amended, and by-laws do not exist, this this article (article VI) allows the creation of by-laws in favor of amending the main body of the constitution (this document).

	It is further strictly discouraged that any amendments or by-laws restricting free commerce and creativity be created. New amendments or by-laws should not conflict with the above stated purpose of the club.

	\section{Article VII - Method of Dissolution of Organization}

	Should the club be forced to dissolve itself, any and all assets should be put toward the club's debt (if any). This includes university owned equipment or university funds.

	The remaining assets (hardware, operating funds, etc) should be donated to a free software organization, like the fsf. 5 The organization to which the assets are donated must be determined at the time of dissolution by the Benevolent Dictator, with approval of the supervisor.

	\begin{center}
		\includegraphics [height=2.5em] {cc-0.png}
	\end{center}

\end{document}
